% ── ABSTRACT (~200 words) ─────────────────────────────────────────────────────
% No mathematics required. State the two-wings principle and structural
% isomorphism finding. Accessible to Bahá'í scholars.

The Bahá'í principle that science and religion are two wings of one truth is
foundational to Bahá'í teaching.
This paper tests that principle against the deepest findings of modern physics:
Bell's theorem~\citep{Bell1964}, relational quantum mechanics~\citep{Rovelli1996},
and the 2022 Nobel Prize in Physics~\citep{Zeilinger2023Nobel}.

The finding is that the Bahá'í metaphysics of \mahabbat (relational love as the
creative foundation of existence) is structurally isomorphic to relational quantum
mechanics and Whitehead's process philosophy~\citep{Whitehead1929} --- not
analogically, but in the precise sense of sharing the same logical architecture:
existence constituted by relationship, properties emerging through relational
interaction, and the force sustaining relationship identified as love.

The Bahá'í writings articulate this structure forty years before \citeauthor{Whitehead1929}
named it philosophically and a century before physics confirmed it experimentally.
This chronological priority is not claimed as proof of supernatural authority;
it is documented as evidence of independent convergence --- the two-wings principle
confirmed at the level of structure, not metaphor.
